% ============================================================
% Chapter 2 — Background and Related Work
% ============================================================
\chapter{Background and Related Work}
\label{ch:background}

\roadmap{This chapter reviews the five strands of literature on which the
platform rests: the lineage of affect-aware intelligent tutoring systems;
the state of the art in webcam-based engagement and affect detection; the
evidence base for the four learning-strategy interventions; the SEEV
tradition of attention-allocation modeling that motivates the AOI layer;
and the executive-function, self-regulated-learning, and motivation
frameworks that ground the text-mining constructs. Each section closes by
stating what the present system takes from that strand and where it
deliberately departs.}

\section{Affect-aware intelligent tutoring systems}
\label{sec:bg-its}

The proposition that tutoring systems should perceive and respond to
learner emotion dates to \citet{picard1997}, who argued that affect is not
noise around cognition but part of its machinery, and to
\citet{kort2001}, who proposed a cyclical model in which learners orbit
through affective quadrants as understanding is constructed and repaired.
The empirical program that followed was anchored by AutoTutor
\citep{graesser2004autotutor}, a natural-language tutor whose dialogues
provided the corpus in which the learning-centered states---engagement,
confusion, frustration, boredom---were first systematically observed
\citep{craig2004}. Two descendants matter here. \emph{Affective AutoTutor}
\citep{dmello2012affective} closed the loop: it detected affect from
dialogue, posture, and facial features and altered its pedagogical and
motivational moves in response, producing learning benefits for
low-domain-knowledge students. \emph{GazeTutor} \citep{dmello2012gazetutor}
did the same with gaze alone, redirecting attention when the eyes wandered
from the display, and demonstrated that even coarse gaze signals can carry
pedagogically actionable information. In parallel, MetaTutor
\citep{azevedo2010metatutor} instrumented self-regulated learning
processes---goal setting, monitoring, strategy use---with multichannel
trace data, establishing the paradigm of the tutoring system as a
\emph{measurement instrument} as much as an instructional agent.
\todo{The proposal's literature review also cited MITERS; add the full
citation from the proposal's reference list and one sentence on its
relevance---the reference could not be reconstructed from the evidence
pack with confidence, and no entry has been fabricated.}

Two lessons from this lineage shaped the present design. First, the
transition structure among affective states is diagnostic: confusion is
productive when resolved and corrosive when it persists
\citep{dmello2012dynamics, dmello2014confusion}, and boredom is the state
most reliably associated with poor outcomes \citep{baker2010better}. A
platform should therefore capture state \emph{sequences} at a temporal
resolution fine enough to observe transitions, not merely aggregate
incidence. Second, closed-loop adaptation is powerful but methodologically
expensive: once the tutor responds to detected affect, the detector's
errors contaminate the pedagogy and the pedagogy contaminates the
measurement. The present system departs from Affective AutoTutor precisely
here: sensing is decoupled from intervention by design
(\cref{sec:interventions-autotrigger}), so that detection validity can be
established observationally before any closed loop is attempted.

\section{Webcam-based engagement and affect detection}
\label{sec:bg-webcam}

The feasibility of estimating learning-centered states from consumer
video is now well established, with equally well-established ceilings.
\citet{whitehill2014} showed that engagement judgments made from faces
alone are reliable across human labelers and learnable by machine
classifiers, framing engagement as something faces \emph{express} rather
than something only behavior reveals. \citet{grafsgaard2013} tied specific
facial action units from automated FACS coding \citep{ekman1978} to
engagement and frustration in tutoring dialogues, licensing the use of AU
intensities---rather than raw pixels or categorical emotions alone---as
interpretable features. \citet{bosch2015, bosch2016} moved detection
``into the wild'' of real classrooms with off-the-shelf webcams,
demonstrating above-chance detection of boredom, confusion, and engaged
concentration under uncontrolled lighting, pose, and occlusion, while
reporting the honest costs: substantial fractions of frames without a
usable face, and accuracies far below laboratory figures.
\citet{monkaresi2017} added video-estimated heart rate to facial features
for engagement detection, and the meta-analysis of
\citet{dmello2015review} quantified the general pattern: multimodal
detectors outperform their best unimodal component by roughly ten percent
on average---a real but modest gain that argues for multimodality as
robustness rather than as a silver bullet. On the attention side,
\citet{bixler2016} detected mind wandering from gaze during reading, and
the BROMP observational protocol \citep{ocumpaugh2015} standardized the
human-coded on-task/off-task and affect labels against which such
detectors are conventionally validated.

The present platform positions itself in this tradition with three
specific commitments. Its facial channel separates \emph{categorical
emotion probabilities} (eight AffectNet classes
\citep{mollahosseini2019affectnet}, produced by OpenFace~3.0
\citep{openface3}, a successor to the OpenFace toolkits
\citep{baltrusaitis2018openface}) from \emph{AU intensities} (eighteen
units produced by Py-Feat \citep{cheong2023pyfeat}), because the two
feature families support different analyses and have different validity
arguments. Its gaze channel uses WebGazer
\citep{papoutsaki2016webgazer}, accepting region-level rather than
fixation-level precision. And, lacking the human BROMP-style live coding
that \citet{bosch2015} used as ground truth, it substitutes two weaker but
attainable standards: periodic learner self-report, and cross-modal
convergent validation between the facial and gaze channels
\citep{campbell1959}---with the retrospective coding surface built so
that human labeling remains possible in future work
(\cref{ch:sensing,ch:future}).

\section{Evidence for the four learning-strategy interventions}
\label{sec:bg-strategies}

The platform's pedagogy layer operationalizes four strategies chosen from
the effectiveness review of \citet{dunlosky2013}, which graded ten
learning techniques by the breadth and robustness of their evidence.

\textbf{Practice testing.} Retrieval practice is one of the two techniques
\citet{dunlosky2013} rated of high utility. The canonical demonstration is
\citet{roediger2006}: tested learners retained prose passages better at a
week's delay than learners who restudied, despite restudy ``feeling''
more effective---an instance of desirable difficulty
\citep{bjork1994}. The platform's practice-testing intervention generates
mixed multiple-choice and short-answer items over the learner's current
material (\cref{sec:interventions-pt}).

\textbf{Distributed practice.} The second high-utility technique.
The meta-analysis of \citet{cepeda2006} across 254 studies found spaced
practice reliably superior to massed practice for verbal recall. The
platform implements spacing through the SM-2 scheduling algorithm of
\citet{wozniak1994}, whose ease-factor update and expanding-interval rules
are quoted verbatim from the implementation in
\cref{sec:interventions-dp}.

\textbf{Interrogative elaboration.} Elaborative interrogation---prompting
learners to generate explanations for \emph{why} stated facts hold---has
moderate-utility evidence \citep{dunlosky2013}, with the generation effect
demonstrated by \citet{pressley1987}: self-generated elaborations aid
memory more than provided ones, and precision of elaboration matters. The
platform's implementation generates ``why/how'' question suggestions and
hosts an explanatory dialogue rather than a quiz
(\cref{sec:interventions-ie}).

\textbf{Stepwise learning (scaffolding).} The stepwise intervention
descends from the scaffolding construct of \citet{wood1976} and its
Vygotskian grounding in the zone of proximal development
\citep{vygotsky1978}: complex material is decomposed into a graded
sequence of steps, each with a comprehension check, with support faded as
competence is demonstrated \citep{vandepol2010}. Decomposition also has a
cognitive-load rationale \citep{sweller1988}: sequencing bounds the
intrinsic load presented at any moment.

Two framing points matter for the evaluation. First, the evidence base for
these strategies comes overwhelmingly from experimenter-controlled
administration; the platform instead makes them \emph{learner-initiated},
which converts the research question from ``does the strategy work?'' to
``do learners elect to use it, and does elected use associate with
learning?''---a self-regulation question as much as a cognitive one.
Second, feedback within the interventions follows the model of
\citet{hattie2007}, whose task/process/self-regulation/self levels
structure the platform's assessment feedback (\cref{ch:modes}).

\section{Attention allocation: SEEV and its proper scope}
\label{sec:bg-seev}

The AOI layer borrows from the SEEV model of visual attention allocation
\citep{wickens2003, wickens2008}, developed to predict where skilled
operators (originally pilots) direct gaze across displays. SEEV holds that
the probability of attending to an area is increased by bottom-up
\emph{salience} and decreased by the \emph{effort} of moving attention,
while top-down \emph{expectancy} (event rate) and \emph{value} (task
importance) pull gaze toward task-relevant regions; well-calibrated
operators allocate attention nearly optimally with respect to expectancy
and value. The model's natural output is a comparison between
\emph{observed} dwell distributions over areas of interest and
\emph{expected} distributions derived from the task.

The platform adopts exactly this observed-versus-expected schema at the
granularity its gaze channel can support: interface panels as AOIs,
task epochs (reading, intervention, dialogue) defining expected
allocations, and a total-variation alignment score per epoch
(\cref{sec:sensing-aoi}). Two scope cautions are inherited deliberately.
SEEV describes \emph{skilled, steady-state} scanning, and its parameters
are normally fitted to expert task analyses; the platform's expected
weights are instead ordinal placeholders supplied by the researcher, so
the resulting \allocscore{} is presented as a relative alignment index,
not a fitted SEEV model. And webcam gaze cannot resolve within-panel
fixations, so all claims are at the panel level. Within those bounds, the
SEEV framing supplies what raw gaze heatmaps do not: a principled,
task-referenced definition of ``attention where it should be.''

\section{Executive function, self-regulated learning, and motivation}
\label{sec:bg-ef}

The text-mining pipeline's nine constructs (\cref{sec:methodology-ef})
are drawn from three adjacent literatures. From the executive-function
tradition, the unity/diversity model of \citet{miyake2000, miyake2012}
identifies updating (working memory), shifting (cognitive flexibility),
and inhibition as separable-but-related control processes; the pipeline's
\emph{working memory} and \emph{cognitive flexibility} constructs are
verbal-marker proxies for the first two, and are flagged in the
implementation itself as low-feasibility, unvalidated-as-EF signals---a
candor this report preserves. From the self-regulated-learning tradition,
the COPES model of \citet{winne1998} and the cyclical-phase account of
\citet{zimmerman2002} treat metacognitive monitoring---the learner's
online evaluation of their own understanding---as the hinge of effective
study; the pipeline's \emph{metacognition} and \emph{metacognitive
monitoring} constructs detect its verbal expressions (``wait, that's
wrong,'' ``I don't understand \emph{this} part''). From the motivation
and emotion tradition, the control-value theory of \citet{pekrun2006}
predicts which achievement emotions (boredom, frustration, anxiety,
enjoyment) arise from which appraisals of control and value, and
self-determination theory \citep{ryan2000} explains why
learner-initiated---autonomy-preserving---interventions may sustain
engagement where system-imposed ones provoke reactance. The pipeline's
\emph{confusion}, \emph{frustration}, \emph{boredom}, and
\emph{engagement/off-task} constructs give these states a linguistic
detection channel that runs in parallel to the facial one, enabling the
cross-channel comparisons described in \cref{ch:methodology}.

\section{Summary}

Prior work establishes that learning-centered affective states are
consequential and detectable from consumer hardware, but the detectors
are noisy, the validation standards are expensive, and the canonical
systems separated sensing research from authentic coursework. The
platform documented in the following chapters responds by embedding a
multimodal, privacy-conscious sensing layer inside a complete LLM
tutoring environment, keeping that layer observational, and building the
replay and coding instruments needed to validate it. The next chapter
accounts for how this shape emerged from the original proposal.
